\documentclass{article}
\usepackage{amsmath,amsthm,amssymb}

\title{Shifted Littlewood--Richardson}
\author{}
\date{}

\newtheorem{conjecture}{Conjecture}
\newtheorem{remark}{Remark}

\begin{document}
\maketitle

\section{Reading words}

This item records the shifted Littlewood--Richardson conjectures from the
source note on skewing \(GP\) and \(GQ\) functions.  The objects are shifted
set-valued tableaux, in either the \(P\)-version or the \(Q\)-version.

The \emph{backword} of such a tableau is read by rows from top to bottom; within
each row, from right to left; and within each box, primed entries first in
decreasing order, followed by unprimed entries in decreasing order.  The
\emph{forword} is read by rows from bottom to top; within each row, from left
to right; and within each box, unprimed entries first in decreasing order,
followed by primed entries in decreasing order.

\section{Lattice and primed-starting conditions}

The lattice property is defined by testing each adjacent pair \(i,i+1\).  In
the backword, count unprimed \(i\)'s and unprimed \((i+1)\)'s; the latter count
must never exceed the former, and an \((i+1)'\) may not occur when the two
counts are equal.  In the forword, count primed \(i'\)'s and primed
\((i+1)'\)'s; again the latter count must never exceed the former, and an
unprimed \(i\) may not occur when the two counts are equal.

The primed-starting property is used for the \(Q\)-rule.  Read boxes by rows
from bottom to top and left to right within each row.  For each \(i\), the first
box containing \(i\) or \(i'\) must contain \(i'\) but not \(i\).

\section{Conjectural rules}

Let \(\mathcal{P}_{\lambda/\mu}\) be the set of shifted set-valued
\(P\)-tableaux of shape \(\lambda/\mu\) satisfying the lattice property.  Let
\(\mathcal{Q}_{\lambda/\mu}\) be the set of shifted set-valued \(Q\)-tableaux of
shape \(\lambda/\mu\) satisfying both the lattice and primed-starting
properties.

\begin{conjecture}[Skew \(GP\) expansion]
\[
  GP_{\lambda/\mu}
  =
  \sum_{P\in\mathcal{P}_{\lambda/\mu}} GP_{\operatorname{wt}(P)}.
\]
\end{conjecture}

\begin{conjecture}[Skew \(GQ\) expansion]
\[
  GQ_{\lambda/\mu}
  =
  \sum_{Q\in\mathcal{Q}_{\lambda/\mu}} GQ_{\operatorname{wt}(Q)}.
\]
\end{conjecture}

\begin{remark}
The source \(GQ\) script checks an equivalent-looking rule in a related \(GR\)
model.  In the source README, \(GR\) is described as the submodel in which, for
each \(i\), the first \(i\) or \(i'\) in the left-to-right, bottom-to-top
reading word is constrained.  The source notes that the \(GR\) check implies
the corresponding \(GQ\) expansion.
\end{remark}

\section{Computer check}

The script \texttt{code/check\_shifted\_lr.py} is a combined curated port of
the source \(GP\) and \(GQ\) scripts.  It compares the direct homogeneous
monomial expansion of a skew shifted Grothendieck function with the homogeneous
monomial expansion reconstructed from the conjectural rule.

The default run checks both the \(GP\) and \(GQ/GR\) branches for degree \(5\),
shape/skew \([3,1]/[1]\), and \(3\) variables.  During curation the same shape
was also checked in degree \(6\).  These computations are finite evidence and
regression checks; they are not proofs of the conjectures.

\end{document}
